\documentclass[10pt, aspectratio=169]{beamer}
\usepackage{../beamer_theme_408}

\title{408考研零基础 --- 操作系统}
\subtitle{3-14 文件共享与保护}
\author{}
\date{}

\begin{document}


\begin{frame}{本节目标}
    \begin{enumerate}
        \item 理解硬链接与软链接的区别
        \item 掌握文件保护的三种方式
        \item 理解访问控制矩阵与 ACL
        \item 熟悉 UNIX 的 rwx 权限模型
    \end{enumerate}
\end{frame}

\begin{frame}{文件共享概述}
    \textbf{文件共享}：多个用户（进程）使用同一个文件。
    \vspace{0.5em}
    \begin{itemize}
        \item 节省磁盘空间（同一文件不必存多份）
        \item 方便协作（多个用户共同操作同一文件）
    \end{itemize}
    \vspace{0.5em}
    两种实现方式：
    \begin{center}
        \begin{tabular}{|l|l|l|}
            \hline
            \textbf{类型} & \textbf{实现机制} & \textbf{典型系统} \\
            \hline
            基于索引节点 & 目录项指向同一 inode & Unix/Linux \\
            基于符号链接 & 创建特殊链接文件 & 所有系统 \\
            \hline
        \end{tabular}
    \end{center}
\end{frame}

\begin{frame}{硬链接（Hard Link）}
    \textbf{硬链接}：多个目录项指向同一个索引节点（inode）。
    \vspace{0.5em}
    \begin{itemize}
        \item 每个硬链接都与原文件共享同一个 inode
        \item 文件数据只有一份，所有链接共享
        \item inode 中维护\textbf{链接计数（link count）}
    \end{itemize}
    \vspace{0.5em}
    删除规则：
    \[
    \text{链接计数} > 1 \Rightarrow \text{仅删除目录项，计数减 1}
    \]
    \[
    \text{链接计数} = 1 \Rightarrow \text{删除文件数据}
    \]
    \vspace{0.5em}
    \textbf{限制}：
    \begin{itemize}
        \item 不能跨文件系统
        \item 不能链接目录
    \end{itemize}
\end{frame}

\begin{frame}{软链接（符号链接）}
    \textbf{软链接/符号链接}：创建一个 \texttt{.link} 文件，存储目标文件的路径。
    \vspace{0.5em}
    \begin{itemize}
        \item 类似于 Windows 的快捷方式
        \item 链接文件有自己的 inode，与被链接文件不同
        \item 访问时 OS 解析路径，定位目标文件
    \end{itemize}
    \vspace{0.5em}
    \begin{center}
        \begin{tabular}{|l|l|l|}
            \hline
            \textbf{特性} & \textbf{硬链接} & \textbf{软链接} \\
            \hline
            inode & 相同 & 不同 \\
            跨文件系统 & 不可以 & 可以 \\
            链接目录 & 不可以 & 可以 \\
            原文件删除后 & 仍可访问 & 失效（悬空链接） \\
            空间占用 & 无额外空间 & 存路径占用少量空间 \\
            \hline
        \end{tabular}
    \end{center}
\end{frame}

\begin{frame}{文件保护方式}
    三种基本保护方式：
    \vspace{0.5em}
    \begin{description}
        \item[口令保护]
            \begin{itemize}
                \item 文件设置口令，访问时需验证
                \item 简单但安全性一般
            \end{itemize}
        \item[加密保护]
            \begin{itemize}
                \item 文件数据加密存储，需要密钥解密
                \item 安全性高，但加解密有性能开销
            \end{itemize}
        \item[访问控制]
            \begin{itemize}
                \item 系统为每个文件维护访问权限列表
                \item 灵活且精细，现代 OS 的主流方案
            \end{itemize}
    \end{description}
\end{frame}

\begin{frame}{访问控制矩阵与 ACL}
    \textbf{访问控制矩阵}：
    \[
    \begin{pmatrix}
        & \text{File}_1 & \text{File}_2 & \text{File}_3 \\
        \text{User}_1 & \text{RW} & \text{R} & - \\
        \text{User}_2 & \text{R} & \text{RW} & \text{RWX} \\
        \text{User}_3 & - & \text{R} & \text{R} \\
    \end{pmatrix}
    \]
    \vspace{0.5em}
    矩阵稀疏大，实际实现常用两种方式：
    \begin{description}
        \item[ACL（访问控制列表）]
            \begin{itemize}
                \item 每个文件关联一个列表，记录哪些用户有什么权限
                \item 按列存储，查询方便
            \end{itemize}
        \item[能力列表（Capability）]
            \begin{itemize}
                \item 每个用户关联一个列表，记录能访问哪些文件
                \item 按行存储
            \end{itemize}
    \end{description}
\end{frame}

\begin{frame}{UNIX 权限模型}
    UNIX 使用简洁的权限模型：\textbf{rwx}
    \vspace{0.5em}
    \begin{center}
        \begin{tabular}{|c|c|}
            \hline
            \textbf{符号} & \textbf{含义} \\
            \hline
            r & 读（Read）--- 4 \\
            w & 写（Write）--- 2 \\
            x & 执行（eXecute）--- 1 \\
            - & 无权限 --- 0 \\
            \hline
        \end{tabular}
    \end{center}
    \vspace{0.5em}
    三组权限：
    \[
    \underbrace{\text{rwx}}_{\text{Owner}} \;
    \underbrace{\text{r-x}}_{\text{Group}} \;
    \underbrace{\text{r--}}_{\text{Others}}
    \]
    \vspace{0.5em}
    数字表示：\texttt{chmod 754 file}
    \begin{itemize}
        \item Owner: rwx = 7
        \item Group: r-x = 5
        \item Others: r-- = 4
    \end{itemize}
\end{frame}

\begin{frame}{本节总结}
    \begin{enumerate}
        \item 硬链接共享 inode（计数删除），软链接存储路径（快捷方式）
        \item 文件保护方式：口令、加密、访问控制
        \item 访问控制矩阵按列存 ACL（每个文件关联权限列表）
        \item UNIX 权限：rwx 三组（Owner/Group/Others），数字表示
    \end{enumerate}
\end{frame}

\end{document}
